\chapter{Related Work}
\label{ch:literature}

Four bodies of work motivate this thesis. Classical auction theory supplies rational-bidder benchmarks. The algorithmic-collusion literature documents how learning agents depart from those benchmarks. The autobidding literature studies how budget-constrained pacing reshapes welfare. Legal scholarship explains why autonomous algorithmic convergence sits awkwardly in antitrust doctrine. I describe each strand in turn, name a small number of anchor papers that the results in later chapters engage directly, and note where this thesis sits relative to each.

\section{Classical Auction Theory}
\label{sec:lit_classical}

Revenue equivalence under independent private values \citep{Vickrey1961}, the affiliation-based revenue ranking favouring second-price formats \citep{MilgromWeber1982}, optimal-reserve design \citep{Myerson1981}, and the marginal-bidder result that an additional bidder raises revenue more than any optimal reserve \citep{BulowKlemperer1996} supply the rational-bidder benchmarks against which I compare algorithmic outcomes. \crefrange{ch:qlearning}{ch:pacing} test whether these predictions survive when bidders are learning algorithms. The dominance of the number-of-bidders factor in every sub-experiment lines up with \citet{BulowKlemperer1996}, and the negative auction-format effect under contextual bandits (\cref{ch:bandits}) lines up directionally with \citet{MilgromWeber1982}. The format reversal observed under pacing (\cref{ch:pacing}) cannot be read against this classical toolkit, because these models assume unconstrained bidders.

\section{Algorithmic Collusion under Learning Algorithms}
\label{sec:lit_algocollusion}

A growing literature asks whether reinforcement-learning agents converge to supra-competitive outcomes in repeated strategic environments, mostly pricing games. Four anchor papers set the frame for my Q-learning and bandit experiments. \citet{Calvano2020} show that Q-learning agents playing Bertrand pricing games converge to prices above the Nash equilibrium and adopt reward-punishment strategies consistent with the Folk Theorem, which states that sufficiently patient players can sustain cooperative outcomes as equilibria of an infinitely repeated game. Their design fixes the auction format and varies only the algorithm's internal settings (learning rate, discount factor, exploration parameter), so the relative importance of algorithmic choices against market structure is not directly addressed.\footnote{Throughout this thesis I use the term hyperparameters for settings that govern how an algorithm learns, such as learning rates, discount factors, regularisation, and exploration intensity. Hyperparameters are distinct from the parameters that the algorithm itself estimates from data.} \citet{Klein2021} extends this analysis to sequential pricing and shows that the granularity of the action space matters, because fine bid grids enable small retaliatory adjustments that sustain coordination. This motivates my choice of an eleven-point action grid, fine enough to permit granular responses but coarse enough for tabular Q-learning to explore.

\citet{BanchioSkrzypacz2022} apply the same reinforcement-learning lens to auctions rather than pricing games and report that first-price auctions are more vulnerable to bid suppression than second-price auctions under Q-learning. The claim that first-price formats are structurally weaker is the headline policy conclusion of the reinforcement-learning auction literature, and I read my results in \cref{ch:qlearning} against this claim directly. My Q-learning results partially corroborate the direction of their claim, at least under affiliated valuations, though the effect is small relative to the number of bidders. My results under pacing in \cref{ch:pacing} flatly contradict it, since first-price auctions generate higher revenue under both dual pacing and PI (proportional-integral) pacing. The reconciliation I propose is that the finding of \citet{BanchioSkrzypacz2022} is specific to the bidding technology they study and does not generalise to the algorithm families deployed on modern advertising exchanges.

Secondary references include \citet{Calvano2023}, who distinguishes genuine collusion from what they call spurious collusion driven by under-exploration, and \citet{Douglas2024}, who prove analytically that deterministic bandit algorithms always converge to collusive outcomes in repeated games while persistent stochastic exploration prevents this. This result partly motivates the memory-decay factor in my LinUCB design, since memory decay is a mechanism for sustaining persistent randomness. Mechanism-level work such as \citet{BanchioMantegazza2022} and \citet{Bertrand2025} identifies specific routes by which Q-learning agents can reach cooperative outcomes, but I do not attempt to pick between these mechanisms in this thesis. I report the outcome patterns and move on.

Empirical studies of real markets motivate the policy concern but cannot isolate drivers. \citet{Assad2020} document margin increases of 9 to 28 percent following the adoption of algorithmic pricing in German gasoline markets, and \citet{BrownMacKay2023} estimate 10 percent price increases among large online retailers that use pricing algorithms. Neither study can distinguish between algorithmic coordination, improved demand prediction, or simple coincidence with other market changes, which is why simulation-based approaches like this thesis remain necessary.

\section{Autobidding, Pacing, and Platform Welfare}
\label{sec:lit_autobidding}

A parallel literature in algorithmic game theory studies budget-constrained pacing in online advertising auctions. Where the algorithmic-collusion literature asks whether learning agents reach supra-competitive prices, the autobidding literature asks how much welfare is lost when budget-constrained advertisers delegate bidding to automated pacing systems. Four anchor papers frame my pacing experiments. \citet{Balseiro2019Learning} introduce dual-based pacing algorithms that adjust a shadow price on the budget constraint, shading bids down when overspending and raising them when underspending, and prove regret guarantees that scale with the square root of the horizon length.\footnote{Regret is the standard performance metric in online learning. It measures the difference between an algorithm's cumulative reward and the reward it would have earned under the best fixed action chosen with hindsight. Regret of order $\sqrt{T}$ after $T$ rounds means the per-round loss relative to the optimal fixed action shrinks to zero as the horizon grows.} This paper supplies the multiplicative dual update used in Experiment~3a.

\citet{Aggarwal2019} prove that uniform bid-scaling (a rule in which a bidder multiplies every valuation by the same constant before submitting it) is optimal for value-maximising agents in truthful auctions, and they establish a worst-case Price of Anarchy of~2 for second-price auctions under budget constraints.\footnote{The Price of Anarchy is the ratio of socially optimal welfare to the welfare realised under self-interested strategic play. A ratio of~1 indicates perfect efficiency; larger ratios indicate larger losses from decentralised behaviour. \cref{ch:auctions_and_algorithms} gives the formal definition used in this thesis.} This bound is the target against which I compare my empirical Price-of-Anarchy estimates in \cref{ch:pacing}, and every observed value falls below it. \citet{Conitzer2022FPPE} prove that first-price pacing equilibria are unique and computable via the Eisenberg-Gale convex programme, a classical convex optimisation problem that also characterises competitive equilibria in the Fisher market model. Their result gives first-price formats a welfare-theoretic advantage that is absent from second-price formats, and it predicts the direction of the auction-format effect I observe in \cref{ch:pacing}. \citet{Gaitonde2023} generalise the welfare analysis by showing that gradient-based pacing achieves at least half the optimal liquid welfare without requiring convergence to equilibrium, a result that holds for any core auction format. This paper supplies the liquid-welfare benchmark used to compute the Price of Anarchy.

Secondary references include \citet{Deng2021TowardsEfficient}, who show that first-price auctions achieve a Price of Anarchy of~1 under return-on-spend constraints, and \citet{FikiorisTardos2023}, who prove that first-price formats remain within a factor of roughly 2.41 of optimum liquid welfare under no-regret pacing (that is, pacing rules whose regret, as defined in the footnote above, grows sub-linearly in the horizon), while second-price formats can be arbitrarily inefficient. \citet{PaesLeme2024} show that autobidding dynamics can exhibit bi-stability and cross-episode drift in simple models, a prediction I test directly in Experiment~3a and do not find in my simulation output. The autobidding literature has yet to measure which design factors matter most for seller revenue. This thesis provides that measurement.

\section{Legal and Regulatory Context}
\label{sec:lit_legal}

Two anchor papers frame the legal question. \citet{EzrachiStucke2017} distinguish four modes of algorithmic participation in collusion, from algorithms as messengers that carry out human-designed cartels to fully autonomous learning agents that converge to collusive outcomes without communication or intent. Their taxonomy is the reference point for debates about whether algorithmic pricing triggers antitrust liability. The most consequential class for this thesis is their fourth category, where the algorithms reach supra-competitive outcomes on their own. My experiments speak to this mode directly. \citet{Hartline2024} proposes regret auditing, in which algorithmic bidders are evaluated against the regret they would experience against the best fixed action rather than against claims about intent. This proposal is methodologically aligned with the black-box framing of this thesis, which treats bidding algorithms as input-output systems and asks which design choices produce the smallest revenue shortfalls.

Both United States antitrust law (Sherman Act \S1) and European competition law (Article~101 TFEU) require evidence of agreement or concerted practice to establish liability for horizontal coordination \citep{Mehra2016}. When algorithms converge to supra-competitive outcomes through independent learning, the evidentiary threshold is not met, and no jurisdiction has sanctioned purely autonomous tacit collusion by independent learning algorithms \citep{OECD2017, OECD2023}. Reform proposals range from per se prohibition of collusion-facilitating designs \citep{Harrington2018} to regret auditing \citep{Hartline2024} and expanded algorithmic inspectability \citep{Gal2019}. \cref{ch:discussion} returns to these debates.

\section{Where This Thesis Sits}
\label{sec:lit_summary}

Each strand of the literature answers a different question. Classical auction theory asks which format maximises revenue for a rational bidder. The algorithmic-collusion literature asks whether learning agents reach supra-competitive prices. The autobidding literature asks how efficient pacing agents are at allocating a fixed budget. Legal scholarship asks when autonomous algorithmic convergence counts as an antitrust violation. The two middle strands reach opposite conclusions about auction format because they study different bidding technologies and never compare them within a single experimental framework. This thesis provides that common framework and tests whether each strand's format recommendation generalises to the other. The short answer, developed in \crefrange{ch:qlearning}{ch:pacing}, is that it does not.
